\scanpage{583}
\mspage{595}
% source-position: upper heading, Remarque and (41)-(42).

\uline{Remark} It is tempting to write the \inserted{first two}
relations (38) in the form
\[
(41)\qquad
\left\{\begin{array}{ll}
(\alpha^{-1}u)(\rho)=\rho
  & \text{i.e.}\quad \alpha^{-1}u\in Z_\rho^{\unsure{\sharp}}\\
(\beta^{-1}u)(\sigma)=\varepsilon_2\sigma
  & \unsure{\text{whence}|\text{i.e.}}\quad \beta^{-1}u\in N'_\sigma\\
\deleted{\unsure{u(\varepsilon_0)}} &
\end{array}\right.
\]
and then to introduce
\[
(42)\qquad
\left\{\begin{array}{ll}
a=\alpha^{-1}u\in Z_\rho^{\unsure{\sharp}},
  & \text{thus }u=\alpha a,\quad \alpha=ua^{-1}\\
b=\beta^{-1}u\in N'_\sigma,
  & \text{thus }u=\beta b,\quad \beta=ub^{-1},
\end{array}\right.
\]
% source-position: middle paragraph and quotient diagram preceding (43).
and the developments of \S\,48 (notably
\uline{\(\overline{\mathrm{XI}}\)}) may lead one to expect that the
quantities \(a,b\), for \(u\in\mathcal{M}[0]\) \unsure{variable},
depend \uline{multiplicatively} on \(u\). We are going to see however
that this is not so. Let indeed
\diagram{\begin{tikzcd}[column sep=large]
\delta:\ \deleted{\unsure{\mathcal{M}}}\ G_{\rho,\sigma}
  \arrow[r,"\sim"] & \mathcal{M}
  \arrow[r] & \gamma_{\rho,\sigma}\simeq\Gamma_{\mathbb Q}
\end{tikzcd}}
be the canonical homomorphism, so that we will have
\[
(43)\qquad
\delta(u)=\delta(b)=\tau_\gamma^{\unsure{\nu_3}}\delta(a),
\]
putting
\[
\tau_\gamma=\delta(\tau),\qquad
\varepsilon_3=(-1)^{\nu_3}\quad
(\nu_3\in\mathbb Z/2\mathbb Z)
\]
% source-position: lower half, (43 bis), (43 ter), and (44).
If \(u'\) corresponds to \(b',a',\) \inserted{\(\nu'_3\)}, we will
likewise have
\[
(43\,\mathrm{bis})\qquad
\delta(u')=\delta(b')=\tau_\gamma^{\nu'_3}\delta(a')
\]
and \unsure{if} \(u''=u'u\) corresponds to \(a'',b'',\nu''_3\),
\[
(43\,\mathrm{ter})\qquad
\delta(u'u)=\delta(b'')=\tau_\gamma^{\nu''_3}\delta(a''),
\]
In fact, we will have
\[
\begin{array}{rl}
(44)\qquad & b''=b'b,\quad \nu''_2=\nu'_2+\nu_2\quad
 \unsure{\text{and.}|\text{and}}\quad \nu''_3=\nu'_3+\nu_3\\
& (\text{i.e. }\varepsilon''_2=\varepsilon'_2\varepsilon_2
\quad\text{and}\quad\varepsilon''_3=\varepsilon'_3\varepsilon_3)
\end{array}
\]
\mspage{582}

% Printed archive footer outside the manuscript frame (not manuscript):
% 582
% https://grothendieck.umontpellier.fr
% Université de Montpellier, UMR CNRS 5149, Case courrier 051, Place Eugène Bataillon
% 34095 Montpellier cedex 5 - France

\scanpage{584}
% source-position: upper calculations leading to (45).

We have seen
\[
\begin{array}{l}
\deleted{\unsure{\varepsilon_2|\text{and}}}u(\sigma)
 =\varepsilon_2\operatorname{int}(\beta)\sigma\\
u'(\sigma)=\varepsilon'_2\operatorname{int}(\beta')\sigma
\end{array}
\]
whence
\[
\begin{array}{l}
(u'u)(\sigma)=u'(u(\sigma))
 =\varepsilon_2\operatorname{int}(u'(\beta))u'(\sigma)\\
\hfill{}=\varepsilon_2\varepsilon'_2
 \operatorname{int}(\deleted{\unsure{\beta'}}u'(\beta)\unsure{\beta')\sigma}
\end{array}
\]
with \(u'(\beta)\beta'\in\widehat{\Pi}_0\), which by uniqueness proves
\[
(45)\qquad
\beta''=u'(\beta)\beta',\qquad
\varepsilon''_2=\varepsilon'_2\varepsilon_2
\]
% source-position: middle, b-coordinate multiplication and underbraces.
whence we indeed conclude
\[
\begin{array}{ccc}
b''&=&\underbrace{b'}_{\beta'^{-1}u'}\ \underbrace{b}_{\beta^{-1}u}\\
\Vert&&\\
\beta''^{-1}u''=\beta''^{-1}(u'u)&&
\end{array}
\]
\[
\begin{array}{lll}
\text{i.e.}&\beta''^{-1}u'u=\beta'^{-1}u'\beta^{-1}u
 &\text{i.e.}\quad\beta''^{-1}=\beta'^{-1}u'(\beta^{-1})\\
\text{i.e.}&\beta''=u'(\beta)\beta'
 &\text{which is none other than (45).}
\end{array}
\]

% source-position: lower middle, epsilon_3 identity and (46).
The relation \(\varepsilon''_3=\varepsilon'_3\varepsilon_3\) is also
evident, \unsure{for example} from the relation (which contains \(\varepsilon_3\),
and likewise \(\varepsilon'_3,\varepsilon''_3\)
\[
(46)\qquad
u(\rho)\deleted{\text{ is}}\ \mathcal{G}^{+\wedge}\text{-conjugate to }\rho^{\varepsilon_3}.
\]
\deleted{\unsure{when we}}

% source-position: lower third, hypothetical a-coordinate law, (*) and final displays.
It remains therefore to examine the hypothetical relation
\[
a''=a'a
\]
which \deleted{\unsure{is written|would be}} implies
\[
(*)\qquad\delta(a'')=\delta(a')\delta(a)
\]
but comparing (43) (43 bis) (43 ter) we find the relation
\[
\deleted{\unsure{\delta}}\delta(u')\delta(u)
 =\tau_\gamma^{\nu'_3}\delta(a')\tau_\gamma^{\nu_3}\delta(a)
\]
\[
\delta(u'\unsure{u''|u})=\tau_\gamma^{\nu''_3}\delta(a'')
\]
\hfill i.e.
\mspage{583}

% Printed archive footer outside the manuscript frame (not manuscript):
% 583
% https://grothendieck.umontpellier.fr
% Université de Montpellier, UMR CNRS 5149, Case courrier 051, Place Eugène Bataillon
% 34095 Montpellier cedex 5 - France

\scanpage{585}
\deleted{\(\unsure{\tau_\gamma^{\nu''_3}}\)}
\mspage{596}
% source-position: upper cancelled sentence, identities and (47).

\deleted{so that the relation (*) is equivalent to}
\[
\tau_\gamma^{\nu'_3+\nu_3}\delta(a'')
 =\tau_\gamma^{\nu'_3}\delta(a')\tau_\gamma^{\nu_3}\delta(a)
\]
or again
\[
\deleted{\tau_\gamma^{\unsure{\nu_3}}}\delta(a'')
 =\deleted{\unsure{\tau_\gamma}}\inserted{\tau_\gamma^{\nu_3}}
 \delta(a')\tau_\gamma^{\nu_3}\delta(a)
 \deleted{\unsure{\tau_\gamma^{\nu_3}}}
\]
i.e.
\[
(47)\qquad
\delta(a'')=\deleted{\unsure{\operatorname{int}}}
 \tau_\gamma^{\nu_3}(\delta(a'))\cdot\delta(a).
\]
% source-position: upper middle, (48) and its left marginal note.
Thus the relation (*) is equivalent to
\[
(48)\qquad
\begin{array}{lll}
\delta(a')=\tau_\gamma^{\nu_3}(\delta(a'))
 &\text{i.e.}&[\delta(a'),\tau_\gamma^{\nu_3}]=1\quad
 \deleted{\unsure{\text{or}}}\\
 &\text{i.e.}&[\delta(u'),\tau_\gamma^{\nu_3}]=1\quad
 \deleted{\unsure{\text{(for}}}
\end{array}
\]
\margin{since \(\delta(u)=\tau_\gamma^{\nu_3}\delta(a)\)}

% source-position: middle, commutation and heuristic comparison with Gamma_Q.
This shows that if \(\nu_3=1\) i.e. \(\varepsilon_3=-1\), then
the relation (*) is equivalent to the \uline{commutation}
\deleted{\unsure{as}}
\uline{of} \(\delta(u')\) \deleted{\(\unsure{\delta(a')}\)}
\uline{with} \(\tau_\gamma\),
\deleted{\unsure{commutation i.e. to}}.
It is a draconian condition on an element of \(\Gamma_{\mathbb Q}^{\sim}\)
\unsure{that|which}, for an element of \(\Gamma_{\mathbb Q}\)
itself \unsure{in its version}, implies that
\(u'\in\{1,\tau_\gamma\}\) --- and by heuristic arguments, we have seen
that it should even be so for \unsure{every element}
of \(\Gamma_{\mathbb Q}^{\sim}\). It appears that the success of the presentation
of \S\,48 \uline{\(\overline{\mathrm{XI}}\)} is due to the fact that
\unsure{in} \unsure{the situations|the situation}
\unsure{considered|considered} in \unsure{the|the} cases,
the group \(\gamma_{\rho,\sigma}\simeq\widehat{\mathbb Z}^{*}\)
was \uline{commutative}, \inserted{or \(\tau_\gamma\) \uline{central}.}
% source-position: lower third, inclusion and restricted multiplicativity; cancelled left margin.
Outside this case, the quantity
\diagram{\begin{tikzcd}
a=\alpha^{-1}u\in Z_\rho^{\unsure{\sharp}}
 \arrow[r,hook] & G_{\rho,\sigma}
\end{tikzcd}}
depends multiplicatively on the element
\(u\in\deleted{\unsure{Z_{\rho,\sigma}}}\mathcal M^{!}[0]\),
only on condition of restricting oneself to \(u\in\mathcal M^{!}[0]'\).

\margin{\deleted{It even suffices that \(\tau_\gamma\)
be \unsure{centralized|central} by}
\diagram{\begin{tikzcd}
\deleted{\gamma^{!}_{\rho,\sigma}=\operatorname{Ker}(\gamma_{\rho,\sigma}}
 \arrow[r,"\deleted{\unsure{\varepsilon_3}}"'] & {}
\end{tikzcd}}}
\mspage{584}

% Printed archive footer outside the manuscript frame (not manuscript):
% 584
% https://grothendieck.umontpellier.fr
% Université de Montpellier, UMR CNRS 5149, Case courrier 051, Place Eugène Bataillon
% 34095 Montpellier cedex 5 - France

\scanpage{586}
% source-position: upper text and homomorphism diagram.

We are going to try to salvage things, by finding a canonical
\uline{homomorphism}
\diagram{\begin{tikzcd}
\mathcal M^{!}[0] \arrow[r] \arrow[d,equal]
 & \underbrace{Z_{\rho,\sigma}^{!}
 \times_{\gamma_{\rho,\sigma}}\deleted{\unsure{Z}}N_\rho^{*}
 \times_{\gamma_{\rho,\sigma}}N_\sigma^{*}}
 _{S_0\Gamma_{\rho,\sigma}^{!}\underset{\text{index }2}{\subset}
 S_0\Gamma_{\rho,\sigma}} \\
Z_{\rho,\sigma}^{!} &
\end{tikzcd}}

\margin{\uline{N.B.} \unsure{we have}
\diagram{\begin{tikzcd}
S_0\Gamma_{\rho,\sigma}^{!}/L_0^{\Pi}\arrow[r,"\sim"]
 & S_0\Gamma_{\rho,\sigma}/L_0^{\mathcal G}
\end{tikzcd}}
\uuline{then}
\diagram{\begin{tikzcd}
\gamma_{\rho,\sigma}\arrow[r,"\sim"] & \Gamma_{\mathbb Q}^{\sim}.
\end{tikzcd}}}

% source-position: upper quarter, annotated triples (49).
Let us note that the second \unsure{side|set} is formed of the triples
\diagram{\begin{tikzcd}[column sep=small,row sep=small]
(49)\quad\bigl\{\bigl(u,\arrow[d,phantom,"\in"{rotate=90}]
 & a,\arrow[d,phantom,"\in"{rotate=90}]
 & \deleted{\unsure{b}}b\bigr)\mid\arrow[d,phantom,"\in"{rotate=90}]
 & u=fb=f'a,\ \text{with }f,f'\in\mathcal G^{+\wedge}
 \text{ suitable}\bigr\} \\
\deleted{\unsure{Z}}\unsure{\mathcal M}[0]
 & N'_\rho & N'_\sigma &
\end{tikzcd}}
\[
\text{which implies}\quad
\begin{array}{l}
\varepsilon_2(u)=\varepsilon_\sigma(b)\\
\varepsilon_3(u)=\varepsilon_\rho(a)
\end{array}
\]
and for such a triple, we will have
% source-position: upper middle, calculations (49 bis) and definition.
\[
(49\ \mathrm{bis})\quad
\left\{\begin{array}{l}
u(\deleted{\unsure{\rho}}\inserted{\sigma})
 =\operatorname{int}(\unsure{fb})
 \deleted{\unsure{\varepsilon_2(b)|\operatorname{int}(b)}}(\sigma)
 =\operatorname{int}(f)\underbrace{\operatorname{int}(b)(\sigma)}
 _{\varepsilon_2(b)(\sigma)}
 =\deleted{\operatorname{int}(f)\unsure{\varepsilon_2}}
 \inserted{\varepsilon_2(u)\operatorname{int}(f)(\sigma)}\\
u(\rho)=\operatorname{int}(f')(\rho^{\varepsilon_3(u)})
 =\operatorname{int}(f')\varepsilon_3(\tau')(\rho)
 =\operatorname{int}(f'\varepsilon_3(\tau'))(\rho)
\end{array}\right.
\]
where we put
\[
\varepsilon_3(\tau')\overset{\mathrm{dfn}}{=}
\left\{\begin{array}{ll}
1&\text{if }\varepsilon_3=+1\\
\unsure{\tau'}&\text{if }\varepsilon_3=-1
\end{array}\right.
\qquad
\begin{array}{l}
\text{i.e. }\varepsilon_3(\tau')=\sigma^{\nu_3}\\
\text{if }\varepsilon_3=(-1)^{\nu_3}
\end{array}
\]

If \(u\) corresponds to \(\alpha,\beta\) (cf. (38)), the
% source-position: middle, equivalent conjugation formulas after (49 bis).
formulas (49 bis) are written \unsure{also|thus} respectively
\[
\left\{\begin{array}{lll}
\operatorname{int}(f)(\sigma)=\operatorname{int}(\beta)(\sigma)
 &\text{i.e.}& f=\beta\sigma^i\quad(i\in\mathbb Z/4\mathbb Z)\\
\deleted{\operatorname{int}(f')\unsure{\varepsilon_3}(\rho)=}&&\\
\operatorname{int}(f'\varepsilon_3(\tau'))(\rho)
 =\operatorname{int}(\alpha)(\rho)
 &\text{i.e.}&\underbrace{f'\varepsilon_3(\tau')=\alpha\rho^j}
 _{j\in\mathbb Z/6\mathbb Z}
\end{array}\right.
\]

Thus we find

% source-position: lower third, Proposition and (50).
\uline{Proposition}
For a fixed \(u\in Z_{\rho,\sigma}^{!}=\mathcal M[0]
\simeq\unsure{\widetilde{\mathcal M}_{0,3}}[0]\),
the \((u,a,b)\in S_0\Gamma_{\rho,\sigma}\) over \(u\)
are exactly the elements \deleted{of the form}\inserted{of}
\[
(50)\quad
\left\{(u,a,b)\left|\begin{array}{lll}
a=f'^{-1}u&\inserted{\text{with}}\quad
 f'=\alpha\rho^j\varepsilon_3(\tau')&(j\in\mathbb Z/6\mathbb Z)\\
b=f^{-1}u&\text{with }f=\beta\sigma^i&(i\in\mathbb Z/4\mathbb Z)
\end{array}\right.\right\}
\]
\margin{where \(\alpha,\beta,\varepsilon_2,\varepsilon_3\) are
defined for \(u\) \unsure{via} (38)}

\deleted{We will have}\inserted{Let us consider} then the element
% source-position: final paragraph following (50), ending at the bottom of the scan.
\inserted{\((u,a,b)\)} over \(u\)
having the simplest expression in terms of \(\alpha,\beta\),
by taking \(i=\unsure{0|j},\ j=0\), thus corresponding to the choices:
\mspage{585}

% Printed archive footer outside the manuscript frame (not manuscript):
% 585
% https://grothendieck.umontpellier.fr
% Université de Montpellier, UMR CNRS 5149, Case courrier 051, Place Eugène Bataillon
% 34095 Montpellier cedex 5 - France